\section{尖点}
\label{sec:cusps}
% ============================================================================

基本域 $\mathcal{F}$ 向上延伸到无穷远——它的"顶部"越来越窄，
像一个无限高的烟囱。这个"烟囱口"在数学上需要一个名字，
它就是\textbf{尖点}（cusp）。

\textbf{为什么需要尖点？}
回忆模形式 $E_4(\tau)$ 的 $q$-展开 $\sum a_n q^n$，其中 $q = e^{2\pi i\tau}$。
当 $\im(\tau) \to +\infty$ 时，$q \to 0$。
所以"$\tau$ 趋向无穷远"对应"$q$ 趋向 $0$"，而我们要求模形式在 $q = 0$
处有良好的行为（全纯，甚至为零）。
这个"无穷远点"就是最基本的尖点 $\infty$。

\textbf{$\mathbb{P}^1(\Q)$ 是什么？}
除了 $\infty$ 之外，还有其他"方向"可以趋向实轴上的有理点。
$\mathbb{P}^1(\Q)$ 就是 $\Q \cup \{\infty\}$——
所有有理数加上一个无穷远点，可以理解为"有理数的圆周"。

\begin{definition}[扩展上半平面与尖点]
\label{def:cusps}
\index{扩展上半平面}\index{尖点 (cusp)}\index{P1Q@$\mathbb{P}^1(\mathbb{Q})$}
\textbf{扩展上半平面}定义为
\[
  \HH^* = \HH \cup \Q \cup \{\infty\} = \HH \cup \mathbb{P}^1(\Q).
\]
$\SL_2(\Z)$ 的莫比乌斯变换自然延拓到 $\mathbb{P}^1(\Q)$ 上：
对 $\gamma = \begin{pmatrix} a & b \\ c & d \end{pmatrix}$，
\[
  \gamma(\infty) = \frac{a}{c}, \qquad
  \gamma\!\left(-\frac{d}{c}\right) = \infty.
\]
（就是把 $\frac{a\tau + b}{c\tau + d}$ 中分母为零和分子为零的情况补上。）

$\mathbb{P}^1(\Q)$ 中的点称为\textbf{尖点}（cusps）。
\end{definition}

\begin{example}[具体的尖点变换]
\label{ex:cusp-action}
\leavevmode
\begin{itemize}
  \item $T(\infty) = \infty + 1 = \infty$
    ——平移不影响无穷远点。
  \item $S(\infty) = -1/\infty = 0$
    ——反演把 $\infty$ 映到 $0$。
  \item $S(0) = -1/0 = \infty$
    ——反演把 $0$ 映回 $\infty$。
  \item 取 $\gamma = \begin{pmatrix} 3 & 1 \\ 2 & 1 \end{pmatrix}
    \in \SL_2(\Z)$（$3 \cdot 1 - 1 \cdot 2 = 1$），
    则 $\gamma(\infty) = 3/2$。
    所以 $3/2$ 和 $\infty$ 在同一轨道中。
\end{itemize}
\end{example}

\begin{proposition}
\label{prop:cusps-transitive}
$\SL_2(\Z)$ 在 $\mathbb{P}^1(\Q)$ 上的作用是传递的，即所有尖点都等价。
换句话说：$\SL_2(\Z) \backslash \mathbb{P}^1(\Q)$ 只有一个点。
\end{proposition}

\begin{proof}
任取 $a/c \in \Q$（$\gcd(a,c) = 1$），由 Bézout 恒等式，存在 $b, d \in \Z$
使得 $ad - bc = 1$。则
$\gamma = \begin{pmatrix} a & b \\ c & d \end{pmatrix} \in \SL_2(\Z)$
满足 $\gamma(\infty) = a/c$。所以每个有理数都和 $\infty$ 在同一轨道中。
\end{proof}

\textbf{同余子群的尖点。}
对 $\SL_2(\Z)$ 来说只有一个尖点等价类，但对更小的同余子群
（如 $\GammaO(N)$），尖点会分裂成多个等价类。
例如 $\GammaO(2)$ 有 $2$ 个不等价的尖点：$\infty$ 和 $0$
（因为 $S \notin \GammaO(2)$，无法把 $\infty$ 映到 $0$）。
这将在\cref{sec:congruence-subgroups}中进一步讨论。


% ============================================================================
